\documentclass[12pt,a4paper]{article}

\usepackage[utf8]{inputenc}
\usepackage[T1]{fontenc}
\usepackage{kotex}
\usepackage{amsmath,amssymb}
\usepackage{graphicx}
\usepackage{hyperref}
\usepackage[margin=2.5cm]{geometry}
\usepackage{booktabs}
\usepackage{array}
\usepackage{multirow}
\usepackage{xcolor}
\usepackage{enumitem}
\usepackage{colortbl}
\usepackage[round]{natbib}
\bibliographystyle{plainnat}

\definecolor{accent}{RGB}{0,80,160}
\definecolor{reviewer}{RGB}{180,30,30}
\definecolor{strategy}{RGB}{0,120,60}
\definecolor{pass}{RGB}{34,139,34}
\hypersetup{colorlinks=true,linkcolor=accent,urlcolor=accent,citecolor=accent}

\title{SlosimAgent: Natural-Language-Driven AI Agent\\for Autonomous SPH Sloshing Simulation}
\author{Imgyu Kim}
\date{}

\begin{document}
\maketitle

% ════════════════════════════════════════════════════════════════
\section{요약}
% ════════════════════════════════════════════════════════════════

기존 LLM 기반 CFD 에이전트(OpenFOAMGPT, ChatCFD, Foam-Agent, MetaOpenFOAM 등 7종)는
모두 OpenFOAM(격자 기반 FVM)만 지원하며, CFD 전문가의 자동화 도구로 설계되어 있다.
체계적 문헌 조사(1,700편, 8개 DB) 결과,
\textbf{SPH 솔버를 사용하는 LLM 에이전트는 0편}이고,
\textbf{비전문가가 자연어로 슬로싱 시뮬레이션을 수행할 수 있는 도구 역시 0편}임이 확인되었다.

본 연구는 이 공백을 메우기 위해 두 가지를 제안한다:
\begin{enumerate}[nosep]
\item CFD를 모르는 슬로싱 실무자가 자연어로 DualSPHysics GPU 시뮬레이션을
설정$\cdot$실행$\cdot$분석할 수 있는 최초의 도메인 특화 AI 에이전트 \textbf{SlosimAgent}.
\item 에이전트 출력의 물리적 신뢰성을 SPHERIC Test 10 벤치마크로 검증한 결과,
SPH 슬로싱 문헌 전반에 정량 검증 관행이 부재함을 발견하고
이에 대응하는 \textbf{M1--M8 정량 검증 프레임워크}.
\end{enumerate}

10개 E2E 시나리오에서 M-A3 파라미터 정확도 \textbf{82.2\%}(Qwen3-32B)를 달성하고,
SPHERIC Test 10 Water Lateral에서 M2=19.5\%, $r$=0.655의 정량 결과를 보고한다.
2$\times$2 factorial ablation(80 runs)으로 도메인 프롬프트가 도구 호출의
\textbf{필요조건}(prerequisite)임을 발견하며,
32B$\approx$8B ($\Delta$=3.8\%pp, 9/10 동일)로 로컬 SLM 배포 가능성을 입증한다.

% ════════════════════════════════════════════════════════════════
\section{핵심 비교 논문}
% ════════════════════════════════════════════════════════════════

\begin{table}[htbp]
\centering
\caption{핵심 비교 논문 7종 — 본 연구의 위치를 정의하는 논문들}
\label{tab:key-papers}
\small
\begin{tabular}{clp{3.5cm}p{6.5cm}}
\toprule
\# & \textbf{논문} & \textbf{저자} & \textbf{본 연구와의 관계} \\
\midrule
1 & OpenFOAMGPT 2.0 & Pandey et al. & RAG+ReAct, OpenFOAM, 450+ sims --- \textit{가장 큰 CFD agent, 격자 only} \\
2 & ChatCFD & Fan et al. & Structured Reasoning, 315 cases, 82.1\% --- \textit{유일하게 물리적 정확도 보고} \\
3 & Foam-Agent 2.0 & Yue et al. & 6-Agent MCP, 110 cases, 88.2\% --- \textit{MCP 최초 적용, 격자 only} \\
4 & MetaOpenFOAM & Chen et al. & RAG multi-agent, 85\% --- \textit{RAG 패턴의 한계점 노출} \\
5 & English et al. & English et al. & mDBC, SPHERIC T10, \textbf{정량 메트릭 0개} (263인용) --- \textit{SPH 검증의 공백} \\
6 & DualSPHysics v5 & Domínguez et al. & 코드 논문 (494인용) --- \textit{우리 솔버의 기반} \\
7 & Grand Challenges & Vacondio et al. & 5대 Grand Challenge, \textbf{GC5: 정량 검증 요구} --- \textit{우리 M1--M8의 동기} \\
\bottomrule
\end{tabular}
\end{table}

% ════════════════════════════════════════════════════════════════
\section{경쟁 시스템 비교}
% ════════════════════════════════════════════════════════════════

\begin{table}[htbp]
\centering
\caption{경쟁 시스템 상세 비교 --- 8개 축으로 본 SlosimAgent의 차별점}
\label{tab:system-comparison}
\small
\begin{tabular}{lcccccccc}
\toprule
\textbf{시스템} & \textbf{LLM} & \textbf{Solver} & \textbf{Deploy} & \textbf{Target} & \textbf{SPH} & \textbf{Slosh} & \textbf{Validation} & \textbf{Multi-fl.} \\
\midrule
OpenFOAMGPT 2.0 & Claude-3.7 & OF & Cloud & Expert & $\times$ & $\times$ & Run/No-run & $\times$ \\
ChatCFD & DeepSeek-R1 & OF & Cloud & Expert & $\times$ & $\times$ & 68.12\% & $\times$ \\
Foam-Agent 2.0 & Claude 3.5 & OF & Cloud & Expert & $\times$ & $\times$ & Expert rev. & $\times$ \\
MetaOpenFOAM & GPT-4 & OF & Cloud & Expert & $\times$ & $\times$ & Run/No-run & $\times$ \\
NL2FOAM & GPT-4 & OF & Cloud & Expert & $\times$ & $\times$ & None & $\times$ \\
Engineering.ai & Multi-LLM & Multi & Cloud & Expert & $\times$ & $\times$ & None & $\times$ \\
SimLLM & GPT-4 & COMSOL & Cloud & Expert & $\times$ & $\times$ & None & $\times$ \\
\midrule
\rowcolor{accent!8}
\textbf{SlosimAgent} & \textbf{Qwen3-32B} & \textbf{DSPH} & \textbf{Local} & \textbf{Non-exp.} & \checkmark & \checkmark & \textbf{M1--M8} & \checkmark \\
\bottomrule
\end{tabular}
\end{table}

\textit{읽는 법:} 표의 마지막 네 열(SPH, Sloshing, Validation, Multi-fluid)이
본 연구의 positioning이다.
7종 에이전트는 네 열 모두 $\times$ --- SlosimAgent만 \checkmark{}.
이것이 ``기존 시스템과 경쟁''이 아니라 ``기존에 없던 카테고리''를 열었다는 주장의 근거다.

% ════════════════════════════════════════════════════════════════
\section{논문 구조}
% ════════════════════════════════════════════════════════════════

% ————————————————————————
\subsection{Abstract}

슬로싱 실무자(탱크 검사원, 구조 엔지니어)는 CFD 전문 지식 없이 시뮬레이션을 수행할 수 없다.
기존 7종 LLM-CFD 에이전트는 CFD 전문가 대상이며 모두 격자 기반 솔버만 지원한다.
본 연구는 자연어로 DualSPHysics SPH 시뮬레이션을 자율 수행하는 최초의 AI 에이전트
SlosimAgent를 제시하고, SPHERIC Test 10 벤치마크 검증을 통해
에이전트 출력의 물리적 신뢰성을 입증한다.

% ————————————————————————
\subsection{1. Introduction}

\textbf{1.1 Background} ---
슬로싱의 공학적 중요성(해양 LNG, 항공 연료 탱크, 지진 저수조).
SPH가 슬로싱에 적합한 이유: meshless, violent free-surface 포착, 격자 재생성 불필요.
LLM 기반 CFD 에이전트의 부상(7종, 2024--2025).

\textbf{1.2 Problem} --- \textbf{4가지 연구 Gap}

\begin{enumerate}[leftmargin=2em]
\item[\textbf{G0}] \textbf{비전문가 접근성 공백} ---
슬로싱 테크니션은 XML 작성, dp/BC 설정, 후처리에 CFD 전문 지식이 필요.
NL$\rightarrow$도구 패턴이 포토닉스(PIC 설계), 방사선치료, 데이터분석에는 존재하나
SPH 슬로싱에서는 \textbf{0편}.

\item[\textbf{G1}] \textbf{SPH 솔버용 LLM 에이전트 = 0편} ---
7종 전부 OpenFOAM(격자 FVM) only.
입자 기반 솔버는 입력 형식(XML attribute-only), 파라미터 체계(dp, h, CFL),
경계 처리(DBC/mDBC)가 근본적으로 다르므로 OpenFOAM 에이전트를 ``SPH로 확장''하는 것은 불가.

\item[\textbf{G2}] \textbf{로컬 오픈웨이트 SLM = 0편} ---
7종 전부 Claude/GPT-4/DeepSeek 클라우드 API 의존.
엔터프라이즈 환경(조선소 망분리, 군수 시설, 원자력)에서 외부 API 호출은
보안 정책상 불가 --- 로컬 GPU 1장으로 동작하는 에이전트가 필수.

{\color{strategy}\textit{[전략] Reviewer가 ``왜 GPT-4 안 비교?''를 물을 것.
답변: scope는 ``로컬 SLM이 가능한가?''이지 ``Cloud보다 나은가?''가 아님.
Cloud 비교는 future work로 명시적으로 남긴다.}}

\item[\textbf{G3}] \textbf{MCP 기반 과학계산 통합 = 0편} ---
기존 에이전트는 커스텀 래퍼, 쉘 스크립트, if-else 분기로 도구를 호출.
Model Context Protocol(Anthropic, 2024)이 등장했으나
과학계산 솔버에 적용한 사례는 \textbf{0편}.
본 연구는 MCP 표준으로 DualSPHysics를 LLM이 tool로써
제어$\cdot$조작$\cdot$수행할 수 있도록 설계한다.
\end{enumerate}

\textbf{1.3 Contribution}

\begin{enumerate}[nosep]
\item CFD 비전문가가 자연어로 SPH 슬로싱 시뮬레이션을 수행할 수 있는 \textbf{최초의 AI 에이전트} (G0 해소)
\item DualSPHysics GPU 솔버를 통합한 \textbf{최초의 LLM 에이전트}: SPH + 로컬 SLM + MCP (G1+G2+G3 해소)
\item SPHERIC Test 10 벤치마크를 통한 에이전트 출력의 \textbf{정량적 물리 검증} + M1--M8 프레임워크
\item 2$\times$2 factorial ablation(80 runs)으로 \textbf{도구 아키텍처의 계층적 의존 관계}를 발견: 도구 커버리지가 성능 병목
\end{enumerate}

% ————————————————————————
\subsection{2. Related Work}

\textbf{2.1 LLM-Based CFD Agents} ---
OpenFOAMGPT 2.0, ChatCFD, Foam-Agent 2.0, MetaOpenFOAM, NL2FOAM, Engineering.ai, SimLLM.
공통점: OpenFOAM only, Cloud API, CFD expert target.
ChatCFD만 물리적 정확도(68.12\%) 측정 --- 나머지는 ``실행 성공/실패''만 보고.

\textit{논증:} 이 7종은 ``전문가를 더 빠르게(workflow acceleration)'' 만드는 도구이지,
``비전문가를 가능하게(capability enabling)'' 만드는 도구가 아니다.
SlosimAgent는 후자를 목표로 한다 --- 따라서 직접 비교보다 positioning이 다르다는 주장.

\textbf{2.2 NL Interfaces for Scientific Simulation} ---
타 도메인 NL$\rightarrow$도구 패턴의 성공 사례:
포토닉스 PIC 설계, 방사선치료 계획, low-code 데이터 분석.
이들은 ``전문 도구의 비전문가 접근''이라는 동일한 패러다임 ---
SPH 슬로싱은 이 패러다임이 아직 적용되지 않은 영역.

\textbf{2.3 SPH Sloshing and DualSPHysics} ---
DualSPHysics v5 (Domínguez et al. 2022, 494인용): 코드 기반.
English et al. (2022, 263인용): mDBC + SPHERIC T10 ---
그러나 \textbf{정량 메트릭 0개} (시계열 그래프만, 오차율 미보고).
Vacondio et al. (2020): Grand Challenge 5 = ``정량 검증 관행 확립'' 요구 ---
\textbf{우리 M1--M8의 직접적 동기}.

\textbf{2.4 MCP and Tool Integration} ---
Model Context Protocol (Anthropic, 2024): AI 도구 호출의 표준화.
2024년 이전: 에이전트 코드 안에 \texttt{if-else} 분기로 도구 호출을 하드코딩.
2025년 이후: MCP가 등장하면서 ``도구 설명(description)을 LLM이 읽고,
필요할 때 자율적으로 호출''하는 패러다임으로 전환.
과학계산 솔버에 MCP를 적용한 사례는 본 연구가 최초.

% ————————————————————————
\subsection{3. System Design}

\textbf{논리:} 각 서브섹션이 하나의 Gap을 직접 해소하는 구조.
G1$\rightarrow$§3.2, G2$\rightarrow$§3.3, G3$\rightarrow$§3.4.
G0는 §4 전체에서 해소.

\textbf{3.1 Architecture Overview}

\begin{center}
\small
자연어 입력 $\xrightarrow{\text{Qwen3-32B}}$
도메인 프롬프트 + Tool Call $\xrightarrow{\text{MCP 15 tools}}$
DualSPHysics GPU Pipeline $\xrightarrow{\text{Docker CUDA 12.6}}$
AI 분석 + 리포트
\end{center}

5단계 파이프라인:
(1) NL 파싱 $\rightarrow$ (2) XML 생성(xml\_generator) $\rightarrow$
(3) 전처리(GenCase) $\rightarrow$ (4) GPU 시뮬레이션(DualSPHysics5.4) $\rightarrow$
(5) 후처리(PartVTK, MeasureTool) + AI 해석.
\textbf{비전문가는 1단계(자연어)만 제공} --- 나머지 4단계는 에이전트가 자율 수행.

\textbf{3.2 Domain-Specialized Agent (G1 해소)}

DualSPHysics는 OpenFOAM과 근본적으로 다른 입력 체계를 가진다:
\begin{itemize}[nosep]
\item \textbf{XML attribute-only 문법}: \texttt{<gravity x="0" y="0" z="-9.81"/>} --- 텍스트 딕셔너리가 아님
\item \textbf{입자 해상도(dp)}가 격자 크기를 대체 --- dp로부터 h, CFL, 입자 수가 결정됨
\item \textbf{경계 처리}: DBC(Dynamic Boundary Condition) / mDBC --- OpenFOAM의 patch와 전혀 다른 개념
\end{itemize}

따라서 OpenFOAM 에이전트를 ``확장''하는 것이 아니라,
DualSPHysics 전용 도구 15개를 새로 설계해야 한다:
xml\_generator(탱크 치수$\rightarrow$XML), gencase(전처리), solver(GPU 실행),
partvtk/measuretool(후처리), stl\_import(CAD 형상),
baffle\_generator(내부 벽), analysis(AI 물리 해석) 등.

\textbf{슬로싱 도메인 특화 시스템 프롬프트(SloshingCoderPrompt)}:
\begin{itemize}[nosep]
\item 4개 절대 규칙(XML 확장자 금지, 시뮬레이션 디렉토리 구조, Docker 경로)
\item Tool 호출 순서 강제: A(직사각형), B(STL), C(배플) 워크플로우
\item motion\_type 가이드: ``피치/흔들어$\rightarrow$sway'', ``기울어/각도$\rightarrow$pitch''
\item 파라미터 자동 결정: dp=min(tank\_x,y,z)/75, timemax=20/freq, fill\_height=0.5$\times$tank\_z
\item 공진 주파수 공식: $f_n = \frac{1}{2\pi}\sqrt{\frac{g\pi}{L}\tanh\frac{\pi h}{L}}$
\end{itemize}

\textbf{2-Stage Fallback 체계}:
(1단계) LLM이 시스템 프롬프트 규칙으로 파라미터를 추론.
(2단계) xml\_generator가 파라미터 범위를 검증 ---
LLM이 틀려도 도구가 잡아주고, 도구 범위 밖이면 LLM 추론에 의존.
이 이중 구조가 EXP-B에서 ``prompt$\times$tool interaction +20.9\%pp''로 나타난다.

\textbf{3.3 Local SLM Deployment (G2 해소)}

\begin{itemize}[nosep]
\item Qwen3-32B via Ollama --- 소비자 GPU(RTX 4090) 1장으로 동작
\item 클라우드 API 불필요 $\rightarrow$ 조선소 망분리, 군수 시설, 원자력 환경 배포 가능
\item tool calling 능력이 핵심 선정 기준: 8B는 tool call 실패율 높음 $\rightarrow$ 32B 필수
\end{itemize}

{\color{reviewer}\textit{[R] Reviewer 예상 질문: ``왜 LLaMA/Gemma/Mixtral 안 비교?''}}

{\color{strategy}\textit{[전략] 정직하게 한계로 인정.
본 연구의 scope = ``로컬 SLM이 가능한가?''를 Qwen3로 입증.
cross-model 비교는 §6.6 future work에 명시.
단, 32B$\approx$8B ($\Delta$=3.8\%pp, 9/10 동일) 결과가 핵심 방어선:
``모델 크기가 중요하지 않다면, 모델 종류도 도구 아키텍처에 의해 흡수될 가능성이 높다.''}}

\textbf{3.4 MCP-Based Tool Integration (G3 해소)}

\textbf{MCP 이전(2024년 이전)}:
에이전트 코드 안에 \texttt{if-else} 분기로 도구 호출을 하드코딩.
``사용자가 슬로싱을 요청하면 $\rightarrow$ DSPH 솔버 호출'' 식의 결정론적 워크플로우.
도구 추가/변경 시 에이전트 코드 전체를 수정해야 함.

\textbf{MCP 이후(2025년)}:
각 도구에 description(자연어 설명)을 부여하고,
LLM이 상황에 따라 \textbf{자율적으로 도구를 선택$\cdot$호출}.
15개 DualSPHysics 도구 = 15개 description + 15개 파라미터 스키마.
이것이 가능한 이유: DualSPHysics documentation(CUDA API, XML 문법)을
도구 description으로 ``번역''한 것.

구현: Go \texttt{BaseTool} 인터페이스 ---
\texttt{Info()} (도구 설명 반환) + \texttt{Run(ctx, ToolCall) $\rightarrow$ ToolResponse}.
Docker Pipeline: GenCase $\rightarrow$ DualSPHysics5.4 GPU $\rightarrow$ PartVTK $\rightarrow$ MeasureTool.
인프라: NVIDIA RTX 4090, CUDA 12.6, Docker compose.

% ————————————————————————
\subsection{4. Evaluation: 에이전트가 작동하는가? (G0 해소)}

\textit{논증 전략:} §4는 ``비전문가가 자연어로 시뮬레이션을 할 수 있는가?''에 답한다.
§5는 ``그 시뮬레이션 결과를 물리적으로 믿을 수 있는가?''에 답한다.
이 두 질문은 독립적이며, 둘 다 Yes여야 시스템이 유의미하다.

\textbf{4.1 NL$\rightarrow$XML Parameter Fidelity (EXP-A)}

\textit{질문: 에이전트가 자연어를 정확한 시뮬레이션 파라미터로 변환하는가?}

실험 설정: 10개 시나리오(직사각형/원통/STL, 물/오일/LNG, sway/pitch)
$\times$ 2 모델(Qwen3-32B, 8B) $\times$ 3 trials = \textbf{60 runs}, temperature=0.

M-A3 채점: 생성된 XML에서 8개 파라미터를 추출하여 ground truth와 비교
(tank$_{xyz}$, fill\_height, motion\_type, frequency, amplitude, timemax).

\textbf{핵심 결과:}
\begin{itemize}[nosep]
\item v3 (baseline): 32B=69.5\%, 8B=64.5\% ($\Delta$=5.0\%pp)
\item \textbf{v4 (P1+P2 fix): 32B=82.2\%, 8B=78.5\% ($\Delta$=3.8\%pp)}
\item 도구 설계 수정(pitch 운동 지원 + 진폭 단위 변환)만으로 \textbf{+12.8\%pp} 개선
\end{itemize}

\textit{이 결과가 말하는 것:}
모델을 바꾸지 않고 \textbf{도구만 고쳤더니} 12.8\%pp가 올랐다.
이것은 ``더 큰 모델''이 아니라 ``더 넓은 도구 커버리지''가
성능 개선의 핵심 경로임을 실증한다.

\textbf{체계적 오류 패턴}:
\begin{itemize}[nosep]
\item P1: motion\_type 혼동(5/10 시나리오) --- xml\_generator가 sway만 지원 $\rightarrow$ v4에서 pitch 추가
\item P2: amplitude 단위 미변환(5/10) --- degrees$\rightarrow$radians 누락 $\rightarrow$ v4에서 수정
\item P3: timemax 과소추정(6/10) --- 프롬프트 규칙 미준수
\item P4: cylinder geometry 미지원(2/10) --- drawbox만 구현
\item P5: tank\_z/fill\_height 혼동(1/10)
\end{itemize}

P1+P2가 도구 수정으로 해결되었고,
P3+P4를 추가 수정하면 $\sim$88--92\%로 상승할 것으로 추정.
\textbf{잔여 오류는 모두 도구 커버리지의 한계}이지 모델의 한계가 아니다.

\textbf{32B $\approx$ 8B}: 10개 시나리오 중 \textbf{9개 동일}, 1개 예외(S05: parametric pitch 추론).
$\Delta$=3.8\%pp = 아키텍처 효과(56.5\%pp)의 1/15.
\textbf{도구 아키텍처가 모델 크기 분산을 흡수한다.}

Determinism: 32B, 8B 모두 10/10 시나리오에서 $\sigma$=0.0 --- temperature=0 조건에서 완전 재현.

\textbf{4.2 Architecture Ablation (EXP-B)}

\textit{질문: 시스템의 어떤 구성 요소가 성능에 기여하는가?}

2$\times$2 factorial design: Domain Prompt $\times$ DualSPHysics Tools.
10 시나리오 $\times$ 2 모델 = 20 runs/condition, 총 \textbf{80 runs}.

\begin{center}
\begin{tabular}{lcccc}
\toprule
\textbf{Condition} & \textbf{Prompt} & \textbf{Tools} & \textbf{M-A3} & \textbf{해석} \\
\midrule
B0 (Full) & \checkmark & \checkmark & \textbf{67.0\%} & 완전 시스템 \\
B1 ($-$Prompt) & $\times$ & \checkmark & \textbf{0\%} & 도구 있어도 0\% \\
B2 ($-$Tool) & \checkmark & $\times$ & 46.1\% & 프롬프트만으로 절반 \\
B4 (Bare) & $\times$ & $\times$ & \textbf{0\%} & 완전 실패 \\
\bottomrule
\end{tabular}
\end{center}

\textbf{핵심 발견 --- Hierarchical Dependency:}

B1=B4=\textbf{0\%} (40/40 runs 전부). 도메인 프롬프트가 없으면
모델이 DualSPHysics 도구를 \textbf{단 1회도 호출하지 않는다} (glob, edit만 사용).
이것은 표준 factorial에서 발생하지 않는 \textbf{퇴화(degenerate) 패턴}이다.

\textit{해석:} 프롬프트는 도구 호출의 \textbf{필요조건(prerequisite)}이고,
도구는 프롬프트 위에서만 작동하는 \textbf{조건부 정제(conditional refinement)}다.
이것이 ``interaction effect''가 main effect보다 큰 이유다.

2$\times$2 Factorial Decomposition:
\begin{itemize}[nosep]
\item Prompt main effect: \textbf{+56.5\%pp} --- 없으면 모든 것이 0
\item Tool main effect: +10.5\%pp --- B1=0에 의한 희석 포함
\item \textbf{Interaction: +20.9\%pp} --- prompt 존재 시에만 유의미
\end{itemize}

Per-model: 8B에서 tool interaction이 더 큼 (24.2\%pp vs 32B의 17.6\%pp)
$\rightarrow$ \textbf{작은 모델일수록 도구 의존도가 증가한다}.
이것이 32B$\approx$8B의 메커니즘이다:
도구가 모델의 자유도를 제한하여 크기에 따른 분산을 흡수.

\textbf{B2 반증거 (도구 없음 조건의 가치):}
프롬프트만 있는 B2에서 LLM이 직접 XML을 생성했는데,
pitch 시나리오에서 \texttt{mvrotsinu}를 \textbf{정확히 생성}함.
이것은 v3에서 도구(xml\_generator)가 pitch를 sway로 강제 변환한 P1 오류의
반증거이며, ``모델은 알고 있었는데 도구가 막았다''를 입증한다.
$\rightarrow$ v4 도구 수정의 직접적 근거.

% ————————————————————————
\subsection{5. Physics Validation: 물리적으로 믿을 수 있는가?}

\textit{논증:} §4에서 에이전트가 ``작동한다''를 보였다.
그러나 ``작동한다 $\neq$ 물리적으로 정확하다.''
에이전트가 생성한 시뮬레이션이 실제 물리 현상을 재현하는지를
독립적 벤치마크로 검증해야 한다.
이를 위해 SPHERIC Test 10(가장 많이 인용된 슬로싱 벤치마크)과
Rafiee 2011 벤치마크를 사용한다.

\textbf{5.1 SPHERIC Test 10 (EXP-C)}

벤치마크: Botia-Vera et al. (2010), Souto-Iglesias (2011),
실험 102회 반복 --- 슬로싱 SPH 연구의 표준 벤치마크.

M1--M8 정량 메트릭 프레임워크:
Vacondio et al. (2020)의 Grand Challenge 5(``정량 검증 관행 확립'')에 대응하여 제안.
기존 SPH 슬로싱 논문들이 시계열 그래프만 제시하고 오차율을 보고하지 않는 관행에 도전.

\textbf{결과:}
\begin{itemize}[nosep]
\item Water Lateral: \textbf{\textcolor{pass}{PASS}} --- M2=19.5\%, $r$=0.655
\item Oil Lateral: \textbf{\textcolor{pass}{PASS}} --- $r$=0.570, M7=4/4 events captured
  \begin{itemize}[nosep]
  \item Artificial viscosity $\rightarrow$ 86.2\% mean peak error (물리적 부적합)
  \item Laminar+SPS 0.1s $\rightarrow$ 37.7\% (개선)
  \item \textbf{Laminar+SPS 10ms (hi-res) $\rightarrow$ 26.4\%} (Peaks 2--4: 6.5\%, 14.4\%, 16.5\%)
  \item 1,176편 서베이 결과: SPH Oil SPHERIC T10의 \textbf{정량적 에러를 보고한 논문 = 0편}. 본 연구가 최초.
  \end{itemize}
\item Roof: \textcolor{orange}{\textbf{PARTIAL}} --- DBC 한계로 천장 압력 과소
\end{itemize}

Trend convergence: dp=8mm($r$=0.460) $\rightarrow$ 4mm($r$=0.655) $\rightarrow$ 2mm($r$=0.697).
단조 증가하나 수렴 감속 (+42\% $\rightarrow$ +6.4\%).
GCI formal criterion 미충족 --- SPH의 비단조 수렴 특성 (Vacondio et al. 2020).

\textbf{5.2 Rafiee 2011 Benchmark}

\textit{목적:} SPHERIC T10과 독립적인 두 번째 벤치마크로 물리 검증의 robustness 강화.

설정: Tank 1.3m$\times$0.1m$\times$0.9m, Sway $a$=0.1m,
$\omega$=3.1165 rad/s ($f \approx f_{natural}$, near-resonance), Water.
비교 대상: Rafiee \& Thiagarajan (2011) 실험 + Kim \& Park (2023) PoF.

P2 (left wall, mid-height) max peak error:
dp=0.004m DBC (100Hz) = \textbf{4.8\%} |
dp=0.002m DBC = 43.4\%.
역설적으로 저해상도(dp=4mm)가 고해상도(dp=2mm)보다 우수 ---
큰 SPH 커널($h$=8.3mm)이 벽면 근처 입자를 더 넓게 감지하여 압력 포착에 유리.

위상 지연: near-resonance에서 SPH 인공점성이 에너지 축적을 지연시켜
시뮬레이션 피크가 실험보다 1--2주기 늦게 발생.
이것은 SPH 슬로싱 시뮬레이션의 알려진 한계이며,
\textbf{위상 독립적(phase-independent)} 피크 비교가 더 적합한 메트릭임을 시사.

\textbf{5.3 Bridge Experiment: Agent XML $\rightarrow$ Physics}

\textit{목적:} §4(파라미터 정확도)와 §5(물리 검증)를 \textbf{직접 연결}.
``에이전트가 생성한 XML이 실제로 물리적으로 유효한 시뮬레이션을 만드는가?''

S02 agent-generated XML (M-A3=100\%) $\rightarrow$ GenCase(105,465 particles)
$\rightarrow$ DualSPHysics GPU(10s) $\rightarrow$ MeasureTool(5 probes).

\textbf{5/5 physics checks PASS:}
\begin{itemize}[nosep]
\item Hydrostatic: 333.5 Pa vs 323.7 Pa expected (\textbf{err 3.0\%})
\item Anti-phase: $r$(left, right) = $-$0.720 (올바른 반대위상 슬로싱)
\item Frequency: 0.789 Hz vs 0.756 Hz expected (\textbf{err 4.4\%})
\item Wall $>$ Center amplitude: 465 Pa vs 213 Pa (올바른 공간 분포)
\item Stability: 105,465/105,465 particles retained (100\%)
\end{itemize}

{\color{reviewer}\textit{[R] 예상 반론: ``M-A3=100\% 케이스만 검증 = 체리피킹.''}}

{\color{strategy}\textit{[전략] 정직하게 인정하고, §6.5에서 M-A3$<$100\% XML의
물리적 영향을 failure mode 분석으로 보완한다.
``P1 오류가 있어도 슬로싱은 발생한다 --- 다만 올바른 모드가 아닐 뿐.''}}

\textbf{5.4 EXP-D: Autonomous Baffle Optimization}

\textit{목적:} 에이전트가 단순 시뮬레이션을 넘어
\textbf{설계 지원(design assistance)}까지 가능함을 입증.
비전문가가 ``슬로싱을 줄이고 싶다''고 말하면
에이전트가 배플을 자동 설계하고 효과를 정량 검증한다.

설정: 0.6m$\times$0.3m$\times$0.4m 탱크, 50\% 수위, sway $f$=0.95Hz.
Baseline(68,450 fluid) vs Center Baffle($h$=0.14m, $t$=0.016m).

결과: SWL 진폭 감소 \textbf{28.5\%} (Left=28.6\%, Right=28.4\% --- 좌우 대칭).
문헌 기대값 30--50\% 범위 근처 --- 물리적으로 타당한 수준.

\textit{G0+G1 증거:} 비전문가가 ``배플 넣어줘''라고 말하면
에이전트가 STL 생성$\rightarrow$XML 삽입$\rightarrow$시뮬레이션$\rightarrow$SWL 비교까지 자율 수행.

% ————————————————————————
\subsection{6. Discussion}

\textbf{6.1 접근성 도약의 의미와 한계 (G0)}

기존 7종은 ``전문가를 빠르게'' --- SlosimAgent는 ``비전문가를 가능하게.''
비전문가 baseline = 0\% (DualSPHysics XML 수동 작성 불가).
0\%$\rightarrow$82.2\%가 핵심 기여.

{\color{reviewer}\textit{[R] ``82.2\%는 production-ready가 아니다.''}}

{\color{strategy}\textit{[전략] 동의. proof-of-concept 수준임을 인정.
Easy/Medium(70--100\%)에서는 초기 탐색(feasibility study)에 실용적.
Hard(20--51\%)에서는 전문가 검토 필수.
비전문가는 출력 검증 불가 --- ``가능성의 입증''이지 ``신뢰할 수 있는 도구''는 아님.}}

\textbf{6.2 Tool Architecture: 도구 커버리지가 성능을 결정한다}

EXP-A와 EXP-B의 교차 분석이 하나의 결론을 가리킨다:
\textbf{``더 큰 모델''이 아니라 ``더 넓은 도구''가 성능 개선의 핵심 경로.}

근거:
\begin{itemize}[nosep]
\item 모델 스케일링 효과: $\Delta$=3.8\%pp (v4)
\item 도구 수정 효과: +12.8\%pp (P1+P2 fix)
\item 아키텍처 효과: +56.5\%pp (prompt, EXP-B)
\item $\rightarrow$ 도구 $>$ 모델 크기, 아키텍처 $\gg$ 모델 크기
\end{itemize}

B2 반증거: 도구 없이 LLM이 직접 pitch XML을 정확히 생성 ---
``모델은 알고 있었는데 도구가 막았다.''
이것은 ``확신을 가진 오답''(confident wrong answer)의 구조:
temperature=0 + 도구의 P1 편향 $\rightarrow$ 결정론적 동일 오류 반복.
모델의 한계가 아닌 \textbf{도구의 구조적 편향}.

\textbf{6.3 모델 크기 무관성의 함의 (G2)}

32B$\approx$8B (9/10 동일, $\Delta$=3.8\%pp).
도구가 출력 형식을 강제 $\rightarrow$ 모델의 ``자유도'' 제한 $\rightarrow$ 크기 무관.

B2(도구 없음): $\Delta$=11.7\%pp --- 도구가 없으면 격차 확대.
$\rightarrow$ \textbf{도구가 모델 크기 분산을 흡수}하는 메커니즘.

함의: 8B 로컬 배포 가능 $\rightarrow$ 소비자 GPU, 엣지 디바이스, 망분리 환경.

{\color{strategy}\textit{[전략] Cloud LLM 미비교를 한계로 인정.
scope = ``local SLM이 가능한가?''이지 ``cloud보다 나은가?''가 아님.
32B$\approx$8B가 방어선: ``모델 크기가 무관하다면 종류도 흡수될 가능성이 높다.''}}

\textbf{6.4 M1--M8 프레임워크의 의의 (Vacondio GC5 대응)}

English et al. (263인용)이 동일 벤치마크에서 정량 메트릭 0개를 보고한 것은
SPH 커뮤니티의 관행을 반영한다.
Vacondio et al.의 Grand Challenge 5 = ``정량 검증 관행 확립.''

본 연구의 M1--M8은 이 요구에 대한 첫 번째 구체적 대응:
M2(mean peak error), $r$(correlation), M7(event count) 등
\textbf{재현 가능한 정량 메트릭}을 정의하고 실제로 적용.
Oil SPHERIC T10에서 정량적 에러를 보고한 논문 = 0편 (1,176편 서베이) ---
본 연구의 26.4\%가 \textbf{최초 보고}.

\textbf{6.5 Failure Mode 분석}

{\color{reviewer}\textit{[R] ``Bridge가 M-A3=100\%만 검증 --- M-A3$<$100\%는?''}}

P1 오류(pitch$\rightarrow$sway): 두 운동 모두 공진 근처에서 자유표면 진동 유발.
``슬로싱이 발생하는가?'' $\rightarrow$ Yes.
``올바른 모드인가?'' $\rightarrow$ No.
P3 오류(timemax 과소): 시뮬레이션 조기 종료 --- 정상 상태 도달 전 데이터는 물리적으로 유효.

\textbf{해석:} 에이전트의 실패 = ``비물리적 결과''가 아닌 ``부정확한 조건의 물리적 시뮬레이션.''
비전문가 관점에서 초기 탐색(feasibility study) 수준에서 유효.

\textbf{6.6 한계 및 향후 연구}

\begin{itemize}[nosep]
\item \textbf{성능}: M-A3 82.2\%. P3+P4 확장 시 $\sim$92\%
\item \textbf{경계}: DBC only --- mDBC는 drawbox normals 미생성으로 폴백 확인
\item \textbf{형상}: 원통형 미지원 (S07, S09 실패 원인). STL은 S10에서 83\% 달성
\item \textbf{모델}: Qwen3 단일 패밀리. LLaMA/Gemma/Mixtral 교차 검증 + Cloud LLM 비교 필요
\item \textbf{사용자 연구}: 비전문가 피험자 대상 NL 프롬프트 작성 능력 미검증
\item \textbf{Self-verification 부재}: 에이전트가 자신의 출력 오류를 감지하는 메커니즘 없음 --- closed-loop refinement는 향후 과제
\item \textbf{SPH 물리 한계}: near-resonance 위상 지연, DBC 천장 압력 과소
\end{itemize}

% ————————————————————————
\subsection{7. Conclusion}

\begin{enumerate}[nosep]
\item \textbf{최초 비전문가용 SPH 슬로싱 AI 에이전트} (G0) ---
M-A3 82.2\% (10 시나리오, 60+ runs),
Bridge 5/5 physics PASS,
자율 배플 최적화 -28.5\% SWL (EXP-D).

\item \textbf{도구 아키텍처가 성능을 결정한다} (핵심 발견) ---
도구 수정(P1+P2)만으로 +12.8\%pp.
2$\times$2 ablation(80 runs): hierarchical dependency
(프롬프트=필요조건 +56.5\%pp, 도구=조건부 +20.9\%pp).
잔여 도구 한계 해결 시 $\sim$92\%.

\item \textbf{로컬 SLM이 충분하다} (G2) ---
32B$\approx$8B (9/10 동일, $\Delta$=3.8\%pp).
도구가 모델 크기 분산을 흡수.
소비자 GPU 1장으로 배포 가능.

\item \textbf{정량 검증 프레임워크} (G1 + Vacondio GC5) ---
SPHERIC T10 (M2=19.5\%, $r$=0.655),
Oil 피크 오차 26.4\% (최초 정량 보고),
Rafiee 2011 (P2 4.8\%),
M1--M8 프레임워크 제안.
\end{enumerate}

% \bibliography{references}

\end{document}
